\documentclass[12pt]{article}
\usepackage[utf8]{inputenc}
\usepackage[T1]{fontenc}
\usepackage{hyperref}
\title{Nenets Language Resources for K-12 Curriculum Development}
\author{Compiled by Claude Sonnet 5.5 (claude-sonnet-5-5) for Andr\'as Kornai}
\date{}
\begin{document}
\maketitle

\section{Introduction}
Nenets (ISO 639-3 yrk; Glottocode nene1249) is a Samoyedic (Uralic) language of northern Russia with two main varieties, Tundra and Forest Nenets. The 2020 census counted about 24,487 self-identified speakers. It is used as a language of instruction in some schools. Forest Nenets differs considerably from Tundra Nenets and has about 1,000 speakers.

\section{Primary Educational Resources}
\begin{itemize}
\item Yamal regional education reports: Yamal published Nenets, Khanty and Selkup textbooks for grades 1--9. The grade 1--4 Khanty and Nenets textbooks were added to the federal list in December 2018, and literature textbooks for grades 2--4 were planned for 2020. These are Tundra Nenets textbooks. \url{https://sever-press.ru/news/sever-press/yamal-created-everyone-uses-the-regional-department-of-education-is-summing-up-the-results-of-the-year}
\item Forest Nenets materials: researchers at HSE, working with teachers and activists and with Rosneft funding, produced elementary-grade teaching materials, a workbook and a language trainer for Forest Nenets. They are planned for use in boarding schools once approved. \url{https://sever-press.ru/news/sever-press/from-a-year-to-a-decade-new-projects-in-yamal-for-preservation-of-indigenous-languages}
\item News on Yamal rural schools and textbooks: \url{https://sever-press.ru/news/sever-press/robots-textbooks-online-consultations-in-moscow-representatives-of-yamal-will-tell-about-the-rural-school-in-the-arctic}
\end{itemize}
The search results gave no direct link to the textbooks themselves, and no primer was identified.

\section{Linguistic and Reference Materials}
\begin{itemize}
\item Irina Nikolaeva, \emph{A Grammar of Tundra Nenets} (Mouton Grammar Library, 2014), listed on a retail page: \url{https://www.ebay.com/p/205624430}
\item Nenets corpus (UHLCS) at the Language Bank of Finland: Gospel of Luke fragments and sample sentences from Tereshchenko's Nenets-Russian dictionary; access requires a research plan. \url{https://kitwiki.csc.fi/twiki/bin/view/FinCLARIN/KielipankkiAineistotNenetsUHLCS}
\end{itemize}
No current Nenets-English dictionary turned up in the search results.

\section{Religious and Community Resources}
None found in the search results.

\section{Digital and Audio Resources}
\begin{itemize}
\item INEL Nenets Corpus (University of Hamburg): 136 annotated texts in Forest and Tundra Nenets, 61,278 tokens, with glosses and English, German and Russian translations; CC-BY-NC-SA. \url{https://inel.corpora.uni-hamburg.de/NenetsCorpus/}
\item User's Guide to INEL Nenets Corpus (Budzisch, Wagner-Nagy, Arkhipov, 2024): \url{https://ebook.ek.szte.hu/index.php/btk-magyarnyelviirodalmi-intezet/catalog/book/357}
\item INEL Nenets Corpus documentation (PDF): \url{https://s3-uhh.lzs.uni-hamburg.de/gwiss-inel-corpora/remote/nenets-1.0/documentation/INEL_Nenets_Corpus.pdf}
\end{itemize}

\section{Curriculum Development Context}
Nenets has the strongest school presence of the four languages: Tundra Nenets textbooks for grades 1--9 exist in Yamal, and part of them are on the Russian federal list. Gaps remain for Forest Nenets, for upper grades and for materials outside Russian-language instruction. Current availability of the textbooks should be checked with the Yamal education department. The grammar and the INEL corpus could support teacher-training and reference materials.

\end{document}